\section{The exact observation quotient}
\label{sec:v78-information}

Fix a compact connected rectangle $Q$, a point $z_0\in Q$, and three
distinct clocks.  Let $\mathscr W$ consist of nonconstant smooth
functions $W$ with $W<T_1$ on a neighborhood of $Q$.  For a positive
smooth function $a$, write
\[
 \mathcal F(W,[a])_j=
 \frac{a(T_j-W)}{\int_Q a(T_j-W)},\qquad [a]=\{ca:c>0\}.
\]
The brackets denote only multiplication by a constant, not arbitrary
reweighting.  Let $\mathscr L$ be the image of this map.  The group of
positive smooth functions modulo constants acts on $\mathscr L$ by
\begin{equation}\label{eq:v78-reweight}
 (b\cdot f)_j=\frac{b f_j}{\int_Q b f_j}.
\end{equation}
This operation changes the source shape while preserving the clocked
action.  It is defined on laws, not on individual recorded endpoints.

\begin{theorem}[Complete invariant of common reweighting]
\label{thm:v78-quotient}
The map $(W,[a])\mapsto\mathcal F(W,[a])$ is a bijection onto
$\mathscr L$.  The absolute action is a complete invariant of the
reweighting action: two law triples are in the same orbit if and only
if their actions agree.  Consequently
\[
                 \mathscr L/\{\text{common positive reweightings}\}
                         \simeq \mathscr W.
\]
The displayed orbit identification is set-theoretic here.  The quotient
topology and a complete metric are defined in
Proposition~\ref{prop:v79-ambient}, and its smooth action charts in
Theorem~\ref{thm:v79-split}.  At the representative level, with the
denominator margins of Theorem~\ref{thm:v77-clock}, recovery of $W$
and $a/a(z_0)$ is locally Lipschitz in each fixed $C^m$ norm.
\end{theorem}
\begin{proof}
The absolute inverse in Theorem~\ref{thm:v77-clock} recovers $W$
from the three laws, and then recovers
\[
 \frac{a(z)}{a(z_0)}=
 \frac{f_1(z)}{f_1(z_0)}\frac{T_1-W(z_0)}{T_1-W(z)}.
\]
This proves injectivity, and surjectivity holds by the definition of
$\mathscr L$.  Formula~\eqref{eq:v78-reweight} sends
$\mathcal F(W,[a])$ to $\mathcal F(W,[ba])$, so each orbit has one
fixed action.  Conversely, if the actions agree, the multiplier
$b=\widetilde a/a$ carries one triple to the other.  The displayed
formulas and the reciprocal estimates in the clock theorem give the
stability assertion.  No statistical sufficiency or finite-sample
Markov kernel is asserted by this quotient construction.
\end{proof}

\begin{proposition}[Local comparison with canonical data]
\label{prop:v78-lens}
Suppose $W_{st}\ne0$ on a rectangle.  Its three-clock law orbit
determines the marked local canonical relation
\[
 \mathcal C_W=\{((s,-W_s(s,t)),(t,W_t(s,t))):(s,t)\in Q\}
\]
and the absolute value $W(z_0)$.  Conversely, this relation, parametrized
by the same endpoint labels, together with $W(z_0)$ determines the law
orbit.  With the extra nuisance shape $[a]$ it determines the laws
themselves.  The canonical relation alone determines $W$ only up to an
additive constant.
\end{proposition}
\begin{proof}
Recovering $W$ and differentiating gives the first assertion.  Conversely
write the covectors read from the relation as $p(s,t),q(s,t)$.  For
realizable canonical data the one-form $-p\,ds+q\,dt$ is $dW$.
On the simply connected rectangle its integral from $z_0$ is independent
of the path.  Adding $W(z_0)$ recovers $W$ and hence its law orbit.
The forward formula recovers the particular laws if $[a]$ is specified.
Two functions with the same differential on a connected rectangle differ
by a constant; adding any admissible constant leaves this canonical
relation unchanged.  Thus the anchor has an actual mathematical role.
\end{proof}

The momenta in this proposition are canonical covectors in scalar
coordinates, not measured Euclidean angles.  Once a boundary embedding
is recovered, its speed converts them to the usual tangential components
of unit velocity.  This is a local comparison on the given branches;
it does not identify the marks or observation domain with those of an
exterior lens problem or of an enriched periodic length spectrum.

\begin{proposition}[Two-clock ambiguity in the weighted-action model]
\label{prop:v78-two-clock}
For two clocks and unknown weight, every nonconstant action with
positive residuals admits nontrivial nearby weighted-action pairs giving
the same two normalized laws.
\end{proposition}
\begin{proof}
Set $r=(T_1-W)/(T_2-W)$ and, for $\lambda>0$ near one, set
\[
 W_\lambda=\frac{T_1-\lambda T_2 r}{1-\lambda r},\qquad
 a_\lambda=a\frac{T_2-W}{T_2-W_\lambda}.
\]
Compactness and strict residual margins make these smooth and admissible
for sufficiently small $|\lambda-1|$.  Direct substitution gives
\[
 a_\lambda(T_2-W_\lambda)=a(T_2-W),\qquad
 a_\lambda(T_1-W_\lambda)=\lambda a(T_1-W).
\]
Normalization erases the factor $\lambda$.  Since $r>0$ and the clocks
are distinct, $W_\lambda\ne W$ when $\lambda\ne1$.
\end{proof}

This proposition quantifies a gauge in the stated weighted-function
space.  It is not a construction of two noncongruent billiards having
the same physical two-clock experiment.  The three-clock sufficiency
theorem does not require such a construction.  Equally, a branch mixture
is not removed by additional clocks in this model: whenever all residuals
are positive,
\[
 \sum_\gamma a_\gamma(T-W_\gamma)
 =\Bigl(\sum_\gamma a_\gamma\Bigr)
       \left(T-\frac{\sum_\gamma a_\gamma W_\gamma}
                       {\sum_\gamma a_\gamma}\right).
\]
The branch-resolution assumption is therefore explicit in the title,
observation convention and main theorem, rather than hidden in notation.
